\documentclass[11pt]{article}
\usepackage[margin=1in]{geometry}
\usepackage{amsmath,amssymb,amsthm,mathtools}
\usepackage{graphicx}
\usepackage{microtype}
\usepackage[hidelinks]{hyperref}

\newtheorem{theorem}{Theorem}
\newtheorem{lemma}{Lemma}
\newcommand{\Bn}{\mathbb B_n}
\newcommand{\A}{A^2_\alpha(\Bn)}
\newcommand{\BMO}{BMO^1_\alpha}
\newcommand{\Lfin}{L^\infty_{\mathrm{fin}}}
\newcommand{\Tfin}{\mathcal T_{L^\infty_{\mathrm{fin}}}}
\newcommand{\cL}{\mathcal L}

\title{A One-Column Counterexample to the Reverse\\
Toeplitz-Algebra Implication}
\author{Candidate solution packet for arXiv:2510.19589\\
\small agent\_lane\_04 (GPT5.6)}
\date{26 August 2026}

\begin{document}
\maketitle

\begin{center}
\fbox{\parbox{0.91\textwidth}{\textbf{Status.}
Candidate counterexample, likely valid, awaiting human review.}}
\end{center}

\section{Source question}

The source is B\'ekoll\`e, D\'efo, and Tchoundja,
\emph{Analysis of Toeplitz Operators with $BMO^1_\alpha$ operator-valued
symbols on $\ell^2$-valued Bergman spaces}, arXiv:2510.19589 (2025)
\cite{BDT}.  Section 13, page 41 asks whether
\[
 b,b^*\in \BMO\bigl(\Bn,\cL(\ell^2,\ell^1\cap\ell^2)\bigr),
 \qquad
 \widetilde b\in L^\infty\bigl(\Bn,\cL(\ell^2)\bigr),
\]
together with
\[
 \lim_{d\to\infty}\|T_bQ_d\|=0,
\]
implies $T_b\in\Tfin$.  Here $Q_d=M_{I_{(d)}}$ is the orthogonal projection
of $\ell^2$ onto $\overline{\operatorname{span}}\{e_j:j>d\}$, acting
pointwise on $A^2_\alpha(\Bn,\ell^2)$.

\begin{figure}[ht]
\centering
\includegraphics[width=\textwidth]{figures/open_problem_crop.png}
\caption{The open question in Section 13, page 41 of the source paper.}
\end{figure}

\section{Proof intuition}

The hypothesis controls only input-coordinate tails.  We therefore choose a
symbol having exactly one nonzero input column; then $T_bQ_d=0$ automatically
for every $d\geq1$.  That column is distributed among infinitely many output
coordinates using pairwise disjoint angular sectors of the ball.  Normalized
Bergman kernels concentrate at boundary points, so each sector produces a
scalar Toeplitz component of norm one.  Hence no output-coordinate tail can
decay.  This is incompatible with $\Tfin$, because its finite-symbol words
have finite row and column support, and norm closure preserves two-sided tail
decay.

\section{Definitions}

Let $P_\alpha$ be the scalar weighted Bergman projection, acting coordinatewise
on $L^2_\alpha(\Bn,\ell^2)$.  For a bounded operator-valued symbol $b$, set
$T_b=P_\alpha M_b$.  The normalized scalar kernel is
\[
 k_a^\alpha(w)=
 \frac{(1-|a|^2)^{(n+1+\alpha)/2}}
 {(1-\langle w,a\rangle)^{n+1+\alpha}}.
\]
Write $P_d=I-Q_d$ for projection onto the first $d$ coordinates.
The algebra $\Tfin$ is the operator-norm closure of finite sums of finite
products of Toeplitz operators whose symbols have only finitely many nonzero
matrix entries, all bounded \cite{BDT}.

We use either standard norm on $\ell^1\cap\ell^2$ (the sum or maximum of the
two norms); the construction is bounded for both.

\section{The counterexample}

\begin{lemma}[A necessary two-sided tail condition]\label{lem:tails}
If $S\in\Tfin$, then
\[
 \|SQ_d\|\longrightarrow0
 \quad\text{and}\quad
 \|Q_dS\|\longrightarrow0.
\]
\end{lemma}

\begin{proof}
If $u\in\Lfin$, choose $N$ larger than every row and column index on which
$u$ has a nonzero matrix entry.  Since $P_\alpha$ acts coordinatewise,
\[
 Q_NT_u=0,\qquad T_uQ_N=0.
\]
For a finite product of such Toeplitz operators, the leftmost factor gives
finite row support and the rightmost factor gives finite column support.
Thus every finite sum of finite products has both tails identically zero once
$d$ is large enough.  If $S$ is an operator-norm limit of such operators, an
$\varepsilon$-approximation immediately gives
$\limsup_d\|Q_dS\|\leq\varepsilon$ and
$\limsup_d\|SQ_d\|\leq\varepsilon$.  Let $\varepsilon\downarrow0$.
\end{proof}

\begin{lemma}[Disjoint sectors with norm-one Toeplitz operators]
There are pairwise disjoint Borel sets $E_1,E_2,\ldots\subset\Bn$ such that
\[
 \|T_{\chi_{E_j}}\|=1\qquad(j\geq1).
\]
\end{lemma}

\begin{proof}
Choose pairwise disjoint relatively open caps $C_j$ of the unit sphere and a
point $\zeta_j\in C_j$.  Let $E_j$ be the radial sector consisting of nonzero
$w\in\Bn$ whose direction $w/|w|$ lies in $C_j$.  These sets are disjoint,
and $E_j$ contains a relative neighborhood of $\zeta_j$ in the closed ball.

The probability measures
\[
 d\mu_{r\zeta_j}(w)=|k_{r\zeta_j}^\alpha(w)|^2\,d\nu_\alpha(w)
\]
converge weakly to the point mass at $\zeta_j$ as $r\uparrow1$.  For
completeness, if $V$ is any relative neighborhood of $\zeta_j$, compactness
gives $c_V>0$ such that
\[
 |1-r\langle w,\zeta_j\rangle|\geq c_V
 \quad(w\in\overline{\Bn}\setminus V)
\]
for all $r$ sufficiently close to one.  Therefore
\[
 \mu_{r\zeta_j}(\Bn\setminus V)
 \leq C_V(1-r^2)^{n+1+\alpha}\longrightarrow0.
\]
Taking $V\cap\Bn\subset E_j$ yields
$\widetilde{\chi_{E_j}}(r\zeta_j)\to1$.  Since $T_{\chi_{E_j}}$ is a positive
contraction and
\[
 \big\langle T_{\chi_{E_j}}k_a^\alpha,k_a^\alpha\big\rangle
 =\widetilde{\chi_{E_j}}(a),
\]
its norm is exactly one.
\end{proof}

\begin{theorem}\label{thm:main}
For every $n\geq1$ and $\alpha>-1$, the question in Section 13 of
arXiv:2510.19589 has a negative answer.  More strongly, there is a bounded
strongly measurable symbol
\[
 b:\Bn\to\cL(\ell^2,\ell^1\cap\ell^2)
\]
whose Hilbert-space adjoint has the same property, such that
\[
 b,b^*\in\BMO\bigl(\Bn,\cL(\ell^2,\ell^1\cap\ell^2)\bigr),
 \qquad
 \|\widetilde b\|_{L^\infty(\cL(\ell^2))}\leq1,
\]
and
\[
 T_bQ_d=0\quad(d\geq1),
 \qquad
 \|Q_dT_b\|=1\quad(d\geq1).
\]
Consequently $T_b\notin\Tfin$.
\end{theorem}

\begin{proof}
For the sectors above, define
\[
 v(w)=\sum_{j\geq1}\chi_{E_j}(w)e_j,
 \qquad
 b(w)x=\langle x,e_1\rangle v(w).
\]
Because the $E_j$ are disjoint, $v(w)$ is either zero or one standard basis
vector.  Hence $b$ is strongly measurable, takes only countably many values,
and is uniformly bounded from $\ell^2$ into $\ell^1\cap\ell^2$.  Its adjoint is
\[
 b(w)^*x=\langle x,v(w)\rangle e_1,
\]
and has the same boundedness and measurability properties.

Every bounded Banach-valued symbol $c$ belongs to the paper's $\BMO$ class:
from its definition,
\[
 \sup_{z\in\Bn}\int_{\Bn}
 \|c\circ\varphi_z(w)-\widetilde c(z)\|\,d\nu_\alpha(w)
 \leq2\|c\|_\infty.
\]
Thus both required $\BMO$ hypotheses hold.

Set
\[
 p_j(z)=\widetilde{\chi_{E_j}}(z)
 =\int_{E_j}|k_z^\alpha(w)|^2\,d\nu_\alpha(w).
\]
Then $p_j(z)\geq0$ and disjointness gives $\sum_jp_j(z)\leq1$.  Moreover,
\[
 \widetilde b(z)x
 =\langle x,e_1\rangle\sum_{j\geq1}p_j(z)e_j.
\]
The output vector has $\ell^2$ norm at most its $\ell^1$ norm, which is at
most one.  Hence $\|\widetilde b(z)\|_{\cL(\ell^2)}\leq1$ uniformly in $z$.

For $d\geq1$, $b(w)Q_d=0$ pointwise, so $T_bQ_d=0$.  On the other hand, for a
scalar $f\in\A$,
\[
 T_b(fe_1)=\bigl(T_{\chi_{E_1}}f,
 T_{\chi_{E_2}}f,\ldots\bigr).
\]
Fix $j>d$.  The $j$th component implies
\[
 \|Q_dT_b\|\geq\|T_{\chi_{E_j}}\|=1.
\]
Since $b$ is pointwise contractive, $\|T_b\|\leq1$, so equality holds.
Lemma~\ref{lem:tails} now shows that $T_b\notin\Tfin$.
\end{proof}

By Theorem 1.19 of the source paper, this $T_b$ is also sufficiently localized.
Thus the example disproves the proposed infinite-dimensional reverse
implication even inside the source's localized Toeplitz class.

\section{Verification, scope, and novelty}

The proof is exact and uses no computation.  The construction satisfies the
question with the stronger identity $T_bQ_d=0$ for every $d\geq1$.  It does
not challenge Proposition 1.16 of the source: that proposition additionally
assumes the missing left-tail condition $\|Q_dT_b\|\to0$, which fails here as
strongly as possible.

A bounded novelty search on 26 August 2026 used the exact arXiv id and title,
the notation $\Tfin$, exact phrases from Section 13, and combinations of
operator-valued BMO, Toeplitz algebra, and reverse implication.  It found the
source paper, mirrors, and summaries, but no later explicit answer.  This is
not an exhaustive bibliographic guarantee.

\textbf{Human-review recommendation.}  Treat Theorem~\ref{thm:main} as a
high-priority candidate counterexample.  In particular, verify the bounded
Banach-valued BMO observation and Lemma~\ref{lem:tails}; neither step hides a
compactness or uniformity assumption.

\begin{thebibliography}{9}
\bibitem{BDT}
D. B\'ekoll\`e, H. O. D\'efo, and E. L. Tchoundja,
\emph{Analysis of Toeplitz Operators with $BMO^1_\alpha$ operator-valued
symbols on $\ell^2$-valued Bergman spaces}, arXiv:2510.19589, 2025.
\end{thebibliography}

\end{document}
